\section{Exact fused--label rigidity at the equality alias}
\label{sec:tpc43-zero-rigidity}

At \(k=0\), the intercept is the fixed integer \(h\).  Hence
\begin{equation}
 (d_\alpha,n_{\alpha,j,0})
 =(d_\alpha,h)=1,
 \label{eq:tpc43-zero-coprimality}
\end{equation}
and every nonzero random atom is indexed by the squarefree product
\begin{equation}
 u_{\alpha,j}
 :=d_\alpha n_{\alpha,j,0}
 =d_\alpha(\ell_\alpha d_\alpha j+h).
 \label{eq:tpc43-fused-zero-label}
\end{equation}

\begin{lemma}[Fused--label injectivity]
\label{lem:tpc43-fused-label-injectivity}
For all sufficiently large \(X\), fixed \(j\in\cJ\), and two physical
rows \(\alpha,\beta\),
\begin{equation}
 u_{\alpha,j}=u_{\beta,j}
 \quad\Longrightarrow\quad
 \alpha=\beta.
 \label{eq:tpc43-fused-label-injective}
\end{equation}
\end{lemma}

\begin{proof}
Equality of the two labels gives
\begin{equation}
 j(\ell_\alpha d_\alpha^2-\ell_\beta d_\beta^2)
 =h(d_\beta-d_\alpha).
 \label{eq:tpc43-fused-collision-equation}
\end{equation}
If the integer in parentheses on the left is nonzero, the left side has
absolute value at least \(j\), whereas
\begin{equation}
 |h(d_\beta-d_\alpha)|
 \le |h|\max_\rho d_\rho<j
 \label{eq:tpc43-fused-size-contradiction}
\end{equation}
by \eqref{eq:tpc43-rigidity-size}.  This is impossible.  Thus the
parenthesized integer vanishes, and
\eqref{eq:tpc43-fused-collision-equation}, together with \(h\ne0\),
forces \(d_\alpha=d_\beta\).  It then forces
\(\ell_\alpha=\ell_\beta\), so the rows coincide.
\end{proof}

The use of \(j\asymp J\) is essential.  Merely knowing
\(|\cJ|=O(J)\) would not prove the strict inequality in
\eqref{eq:tpc43-fused-size-contradiction}.

\begin{theorem}[Exact equality--column diagonalization]
\label{thm:tpc43-exact-zero-diagonalization}
Put
\begin{equation}
 Z_0(\varepsilon)
 :=\|\mathscr T_{0,\varepsilon}^L\|^2
   +\|\mathscr T_{0,\varepsilon}^R\|^2.
 \label{eq:tpc43-zero-random-energy}
\end{equation}
Then
\begin{equation}
 \boxed{\E_\varepsilon Z_0(\varepsilon)=\cD_0.}
 \label{eq:tpc43-exact-zero-expectation}
\end{equation}
Consequently, for every \(\lambda>0\),
\begin{equation}
 \Pp_\varepsilon\bigl(Z_0(\varepsilon)>\lambda\cD_0\bigr)
 \le\lambda^{-1}.
 \label{eq:tpc43-zero-Markov}
\end{equation}
At the inherited endpoint,
\begin{equation}
 \Pp_\varepsilon\bigl(
 Z_0(\varepsilon)>K_B\cD_0
 \bigr)
 \le K_B^{-1}=X^{-1/400+o(1)}.
 \label{eq:tpc43-zero-endpoint-probability}
\end{equation}
Thus all but at most a proportion \(X^{-1/400+o(1)}\) of the
multiplicative prime--sign environments satisfy the random--model
equality benchmark at the same numerical scale as
\eqref{eq:tpc43-physical-zero-target}.  This statement does not select
the physical all--minus environment.
\end{theorem}

\begin{proof}
For fixed \((\gamma,j)\), the nonzero row atoms have distinct Walsh
labels \(u_{\alpha,j}\) by
\cref{lem:tpc43-fused-label-injectivity}.  Equation
\eqref{eq:tpc43-Walsh-orthogonality} therefore deletes every unequal row
pair and leaves
\begin{equation}
 \E_\varepsilon
 |[\mathscr T_{0,\varepsilon}^L]_{\gamma,j}|^2
 =\sum_\alpha
 |a^L_{\alpha,\gamma,j,0}|^2
 \mu^2(n_{\alpha,j,0}).
 \label{eq:tpc43-zero-coordinate-diagonal}
\end{equation}
Sum over the left coordinates and repeat on the right to obtain
\eqref{eq:tpc43-exact-zero-expectation}.  The probability statements
are immediate if \(\cD_0=0\), since then the nonnegative variable
\(Z_0\) vanishes almost surely.  Otherwise they are Markov's inequality,
followed by
\eqref{eq:tpc43-alias-ledger} and the inherited diagonal bound.
\end{proof}
